% Throughout the six sites, people ask variety of questions toward the trashbot in order to understand its purpose on the site. across the 87 interviews we conducted, although the questions they are varied in kinds, they shared a similar sense making process. The questions they asked, Regardless of sites, tend to  interpret the robot in these way: they overall interpretation methods could be generalized through some connected and related questions, which can be summarized by the "What, how, and why" - "What: identifying the robot and its role", "How: understanding operation and participation", "Why: assessing the reason for its presence"

% Participants most often evaluated whether the robot's work was needed, judging whether the job was already being done by the sanitation arrangements they knew. 
%   \begin{quote}
%     \emph{``I mean, just to see any type of effort to clean up
%     the streets is great, because we need that. Especially at
%     this point. We need that. It's essential. Sometimes, you
%     know, a lot of the cans are overflowing or not available
%     or not maintained correctly. So, to see some type of effort
%     mixed with technology, that's pretty cool.''} (LI-P03)
% \end{quote}
%   Participants also asked whether the robot belonged in the place and what would happen to it there, applying the reputation of the place to the robot's survival (theft, stripping) and to its fit. Participants expressed worried about the robots deployed at the open area and suggested it might be a better fit at some protected area
%   \begin{quote}
%     \emph{``Conceptually, I think it's a great idea. Practically,
%     I think you're going to figure out where it works and
%     where it doesn't... You know where it's going
%     to work? Somewhere like a gated community, somewhere like
%     maybe an isolated community.''} (AP-P04)
% \end{quote}

% \begin{quote}
%     \emph{``Is it going to get in pedestrian's way? Like,
%     is it supposed to be, like, in a confined, contained area?
%     Because maybe it would work better inside...''}
%     (LI-P26)
% \end{quote}

% \paragraph{Inferring the robot's role and institutional responsibility}
% Participants interpreted the robot's role through their personal prior experience together with their understanding
% and expectations of the place where they encountered it. They
% connected its presence to institutions and activities they associated
% with that setting, such as sanitation services, university research,
% or policing, to infer what the robot might be doing there and who
% might be responsible for it.

% At Astoria Park, AP-P01 also connected a possible public authority
% behind the robot with concerns about observation:
% % Source: AP-P01.txt, 01:11.
% \begin{quote}
%     \emph{``This is very funny, but then I thought more, I was like maybe it's a little invasive
%     of privacy for. I'm thinking the city has, like, cameras on
%     this, and this is, like, a city thing.''} (AP-P01)
% \end{quote}
% In this account, interpreting the robot as a city initiative
% accompanied a shift from amusement toward concern about privacy.

% At the Roosevelt Island campus caf\'e, MC-P08 identified the robot
% as a student project and contrasted it with a possible corporate
% deployment:
% \begin{quote}
%     \emph{``Yeah, but, I mean, again, in this case, I could tell
%     it's like a student project... I think if it would look more corporate, maybe the
%     feeling would be different, right?''} (MC-P08)
% \end{quote}
% The distinction between a student project and a corporate product
% provided an institutional frame for interpreting the encounter.
% The participant's comparison also connected perceived affiliation
% to the robot's appearance, which associate institute's reputation with the robot's appearance. 

% Institutional responsibility further concerned the purposes for
% which the robot's data could be used. At the Brooklyn Navy Yard,
% when asked about concerns regarding future autonomous robots,
% NY-P09 distinguished police access from monitoring for maintenance:
% \begin{quote}
%     \emph{``So if the police were able to use that data, I don't
%     think that's okay. But for a regular monitoring station to
%     see if the can is full, did it tip over, is it having trouble
%     getting around, is it stuck or something like that, that's
%     fine. That's part of the application.''} (NY-P09)
% \end{quote}
% This account established different expectations for institutions
% that might use the robot's data, without identifying its actual
% operator.

% Together, these accounts show how institutional knowledge and
% expectations supplied resources for interpreting the robot's role.
% Participants situated its possible activities within familiar
% arrangements for maintaining, studying, and monitoring public
% space, while considering the responsibilities associated with
% those activities.

% \paragraph{Inferring participation through the robot's operation}
% TrashBot's movement and visible camera distinguished it from
% an ordinary trash barrel. Participants treated these features
% as clues to how it operated, using them to consider how the
% robot might respond to their actions and what engaging with
% it might involve.

% \emph{Interpreting movement as a response.}
% At the Roosevelt Island caf\'e, a participant in the MC-P01
% interview recalled waving to bring the robot closer, while
% remaining uncertain whether its movement responded to that
% gesture:
% % Source: MC-P01.txt, 02:30--02:38.
% \begin{quote}
%     \emph{``I was confused. I need to wave to him,
%     it would come closer, so I waved, and I don't know if
%     it attracted to my move or if it approached to me
%     by itself.''} (MC-P01)
% \end{quote}
% The participant treated the robot's movement as potentially
% responsive and tried a gesture to coordinate its approach.
% Their uncertainty concerned the relationship between their
% own action and the robot's response: whether waving summoned
% it or its approach occurred independently.

% \emph{Interpreting the camera as observation.}
% At Starlight Park, the camera prompted consideration of what
% being near the robot might entail. A participant in the
% SP-P01 group interview anticipated how other neighborhood
% residents might interpret it:
% % Source: SP-P01.txt, 03:48--03:56.
% \begin{quote}
%     \emph{``And I think that's why it needs to be a little
%     more people friendly. Because in this neighborhood,
%     I feel like a camera is going to make folks feel like,
%     who's watching? You know? Is this the police?''} (SP-P01)
% \end{quote}
% This interpretation connected the visible camera with
% possible observation by an outside authority. The participant
% linked that anticipated reading to the need for a more
% approachable robot, suggesting that understanding participation
% also involved considering who might observe an encounter
% and what that observation could mean.

% In Little Italy, LI-P07 thought the bins served the nearby
% stores and merchants and interpreted the camera as monitoring unwanted
% dumping, raising uncertainty about whether passersby could
% use them:
% \begin{quote}
%     \emph{``Yeah, I thought the camera was, like---I thought
%     the garbage cans were for this place, and I thought the
%     camera was---they wanted to see, make sure nobody was
%     throwing their garbage.''} (LI-P07)
% \end{quote}
% This interpretation connected the camera's apparent function
% with uncertainty about permission to participate. If the
% camera monitored unwanted disposal, placing waste in the
% robot could constitute misuse rather than the intended
% interaction. The participant therefore questioned whether
% passersby were allowed to deposit their rubbish. Here,
% interpreting the camera involved working out the rules
% governing access to the robot's service.

% At the Manhattan caf\'e, two participants in the MC-P02 interview proposed different explanations of how the camera supported the robot's operation. One participant explained:
% \begin{quote}
%     \emph{``I thought it was, like, trying to, you know, kind of,
%     like, turn the camera to, like, see where the trash is,
%     so to pick it up.''} (MC-P02)
% \end{quote}
% The other participant offered a different interpretation:
% \begin{quote}
%     \emph{``No, I think it was looking at the people to see if
%     they have trash to throw away.''} (MC-P02)
% \end{quote}
% These accounts assigned different targets to the camera:
% discarded objects or people with waste to deposit. They also
% implied different forms of participation, depending on whether
% the robot collected waste itself or located people who would
% place waste inside it.

% Questions about participation also concerned who would approach
% whom. At the Brooklyn Navy Yard, NY-P12 contrasted familiar,
% stationary trash cans with the possibility of a robot that
% actively sought out users:
% \begin{quote}
%     \emph{``Well, you gotta think about it. Trash cans are
%     funny because usually, like, I know where a trash can is,
%     because I walk here a lot. But is the intent to have,
%     like, the trash can follow, go near people so that if
%     they have trash, they can throw it away?''} (NY-P12)
% \end{quote}
% The participant sought to establish how access to the service
% would be coordinated: whether people needed to find the robot
% or could expect it to approach them.

% Robot-initiated movement also raised questions about the limits
% of an approach. Recalling their initial reaction at Astoria Park,
% AP-P01 described uncertainty about how close the robot would come:
% \begin{quote}
%     \emph{``I think that my initial reaction also was,
%     I was like trying to get it away. I was like, how close
%     is it going to get to me?''} (AP-P01)
% \end{quote}
% This account shows that understanding participation also involved
% anticipating the robot's behavior when a person wanted distance.
% Across these accounts, participants sought to understand how the
% robot's sensing and movement related to their own actions,
% including depositing waste, waiting for an approach, or trying
% to keep the robot away.

% Participants attended to the robot's appearance when identifying what it was. They are asking what am I encountering? Specifically, what kind of object, intended service, and deployment? In a two-person interview MC-P02, a participant recalled,
% \emph{``Yeah, we thought it was a trash can.''} A later comment in
% the same interview emphasized its familiar form:
% \begin{quote}
%     \emph{``But it's just like a regular trash can, it's good,
%     that's the point.''} (MC-P02)
% \end{quote}

% However, the two participants did not express the same understanding
% of what the robot would do. One raised a question about its
% capabilities:
% \begin{quote}
%     \emph{``And my first thought was like, where's the arm?''...I thought it was going to pick up trash,
%     and I was like, how does it pick up trash?''} (MC-P02)
% \end{quote}

% A question about the intended service also arose at the Brooklyn
% Navy Yard. NY-P12 asked whether the robot was meant to get closer to people:
% \begin{quote}
%     \emph{``Is the intent to have, like, the trash can follow,
%     go near people so that if they have trash, they can throw
%     it away?''} (NY-P12)
% \end{quote}
% Echoing the questions raised in MC-P02, this question sought to
% establish what adding movement to a trash can was intended
% to accomplish.

% The question that mattered most was who was behind the robot and who was responsible for it. People are asking, How does it work, how am I involved, and what people or processes support its operation? Candidate answers came from the setting: the city, a university or research project, the police or ICE, a business association.
  
%   What the camera meant depended on what the place is used for and on who was assumed to be watching.

%    In Little Italy,
% LI-P33 asked, \emph{``Your company is doing this?''} (LI-P33).
% After learning of the team's affiliation, LI-P33 checked
% both the institution and the nature of the activity:
% \begin{quote}
%     \emph{``So this is an experiment?''}
% \end{quote}
% At Astoria Park, AP-P01 recalled an encounter that prompted
% curiosity about both what was happening and who might be watching:
% \begin{quote}
%     \emph{``I literally took a video, I was like, what is going
%     on right now? Oh my gosh, and you caught me when I was
%     smoking... I was like, who is
%     watching behind the scene?''} (AP-P01)
% \end{quote}
